\documentclass[11pt]{article}
% Diplomatic transcription macros (12 Nov 2026). Typed text is plain; the handwritten layer is blue,
% deletions red, typewriter strike-outs grey, doubtful readings in brackets with a question mark.
\usepackage[utf8]{inputenc}
\usepackage[T1]{fontenc}
\usepackage{lmodern}
\usepackage[a4paper,margin=2.2cm]{geometry}
\usepackage{amsmath,amssymb}
\usepackage{xcolor}
\usepackage{tikz-cd}
\usepackage{environ}
\setlength{\parindent}{0pt}
\setlength{\parskip}{4pt}
\title{Nouveau EGA V (cote 25)\\[4pt]\large Diplomatic transcription from the Montpellier scans}
\author{AI-transcribed by Claude Opus 5.5}
\date{12 November 2026}
\newcommand{\scanpage}[1]{\par\bigskip\noindent{\small\sffamily[scan #1]}\par}
\newcommand{\mspage}[1]{{\small\sffamily[p.\,#1]}}
\newcommand{\hand}[1]{\textcolor{blue}{#1}}
\let\over\relax
\newcommand{\over}[2]{\textcolor{gray}{#1}\textcolor{blue}{\ensuremath{^{\!\nearrow}}#2}}
\newcommand{\inserted}[2][]{\textcolor{blue}{$\langle$#2$\rangle$}}
\newcommand{\margin}[1]{\textcolor{blue}{\{\textsf{marg.:} #1\}}}
\newcommand{\deleted}[1]{\textcolor{red}{[\textsf{del.:} #1]}}
\newcommand{\xout}[1]{\textcolor{gray}{[\textsf{x:} #1]}}
\newcommand{\uline}[1]{\underline{#1}}
\newcommand{\uuline}[1]{\underline{\underline{#1}}}
\newcommand{\unsure}[1]{[#1?]}
\newcommand{\illegible}[1]{[\textsf{illeg.\ #1}]}
\newcommand{\diagram}[1]{\par{\centering #1\par}}
\newenvironment{diagramblock}{\par\begin{center}}{\end{center}\par}

\begin{document}
\maketitle

% Nouveau EGA V (cote 25): archive cote 25, site page 2 (PDF page 3).
% typescript. Draft compared: none.
% AI-transcribed by Claude Opus 5.5 (claude-opus-5-5, adaptive thinking, effort high), Message Batches API, batch pilot (msgbatch_0189Nhf1xB615YWnztubtLfc), 2026-11-12; not yet checked by eye.
\scanpage{25/2}
\margin{\unsure{Nouveau EGA V}} \mspage{0}

EGA IV 20 \quad \uline{Sections hyperplanes et projections coniques.} (Etat \unsure{000|0})

1. Notations.

2. Section hyperplane générique: propriétés locales.

3. Section hyperplane générique: irréductibilité et connexité géo\unsure{métriques.}

4. Section hypérplane variable: sections ``assez g\unsure{énérales|éńérales}''.

5. Théorèmes du\unsure{type|ttype} Seidenberg.

6. Connexion d'une section hyperplane quelconque.

7. Application a des constructions de sections hyperplanes et multisec-%
tions spéciales.

8. Dimension de l'ensem\xout{m}ble  des hyperplans exceptionnels.

9. Changement d'immersion projective.

10. Pinceaux de sections hyperplanes, et fibrations de variétés éclatées

11. Grassmaniennes.

12. Généralisation de résultats antérieurs à des sections linéaires.

13. Morphismes élémentaires et théorème de M. Artin.

14. Projections coniques.

15. Axiomatisation de certains des résultats précédents.

16. Traductions en langage de systèmes linéaires.

\uline{Commentaires}. Cette rédaction donne pêle mêle  une esquisse détaillée
de l'ensemble des résultats qui devraient figurer dans une rédaction
finale. Pour parvenir à cette dernière, il faudra une réorganisation de
fond en comble  du présent Etat 0, la première tâche étant probablement
l'établissement d'un nouveau plan (dans lequel sans doute  les actuels
nºs 11,12,14,15  viendront beaucoup plus tôt). Je n'ai pas encore rédi\unsure{-}
le nº 16 , qui en principe ne doit offrir aucune difficulté ni
influer en aucune manière sur les nºs précédents, puisqu'il s'agit d'u\unsure{n}

\margin{2}

\clearpage
% Nouveau EGA V (cote 25): archive cote 25, site page 4 (PDF page 5).
% typescript. Draft compared: none.
% AI-transcribed by Claude Opus 5.5 (claude-opus-5-5, adaptive thinking, effort high), Message Batches API, batch pilot (msgbatch_0189Nhf1xB615YWnztubtLfc), 2026-11-12; not yet checked by eye.
\scanpage{25/4}
IV 20 \quad \uline{Sections hyperplanes et projections \deleted{linéaires}\inserted{coniques}}.

20.1. \quad \uline{Préliminaires \over{de}{\unsure{et}} notations}.

Soient S un \quad préschéma , $\underline{E}$ un Module loc libre de type fini sur S,
$\check{\underline{E}}$ son dual. Nous noterons par P=$\underline{P}(\underline{E})$ le fibré projectif défini par $\underline{E}$,
par $\check{P}$ celui défini par $\check{\underline{E}}$ . Nous\hand{|}appelerons \ $\check{P}$ \ le \uline{schéma des hyperplans}
de P, cette terminologie se justifiant de la façon suivante. \ Soit
$\hand{\xi}$ une section de $\check{P}$ sur S, qui est donc\hand{/}déterminée par un module quotient
inversible L de $\check{\underline{E}}$ . On en conclut un module quotient inversible
$\underline{L}_P$ de$(\check{\underline{E}})_P$ = $(\underline{E}_P)^{\hand{\vee}}$ , d'autre part on dispose d'un module quotient inversibl
$\underline{O}_P(1)$ de \ $\underline{E}_P$ . Passant \ aux duals, on peut envisager \ $\underline{L}_P^{-1}$ resp. $\underline{O}_P(-1)$
comme des sous-modules inversibles (loc facteurs directs) de \ $\underline{E}_P$ resp.
de $(\underline{E}_P)^{\hand{\vee}}$ , \ et l'accouplement \ $\underline{E}_P \otimes \check{\underline{E}}_P \longrightarrow \underline{O}_P$ \ définit \ donc un
accouplement naturel
(x) \[ \underline{O}_P(-1)\otimes\underline{L}_P^{-1} \to \underline{O}_P \]
ou encore un homomorphisme \ transposé
(xx) \[ \underline{O}_P \longrightarrow \underline{O}_P(1)\otimes\underline{L}_P \ = \ \unsure{\underline{L}_P(1)|\underline{D}_P(1)} \ , \]
i.e. une section de $\underline{L}_P(1)$, canoniquement déterminée par $\hand{\xi}$ . Le''diviseur''
de cette section, \ i.e. le sous-schéma\inserted{fermé} $H_{\hand{\xi}}$ de \ P défini par l'Idéal
image de (x), \ est \ appelé l'\uline{hyperplan dans P} \ défini par \ l'élément
$\hand{\xi} \hand{\in} \check{P}(S)$ . \ Il peut encore se décrire en notant que localement sur S \xout{la}
\xout{\unsure{donnée de $\check{\xi}$ équivaut}} \quad $\hand{\xi}$ \ est donné par une section \ $\phi$ de \ $\underline{E}$ , \xout{de} telle
\xout{\illegible{1}} que $\phi(s) \neq 0$ pour tout s \ ( $\phi$ étant déterminé par $\hand{\xi}$ \ modulo
multiplication par une section inversible de $\underline{O}_S$ ); comme $\underline{E}$ = $p_{\hand{*}}(\underline{O}_P(1))$
( p: P$\to$S étant la projection), \ $\phi$ \ peut être considéré comme une section
de \ $\underline{O}_P(1)$, don\over{c}{t} le diviseur n'est autre que $H_{\hand{\xi}}$. \ Bien entendu, si on
\inserted[typed]{\uline{considére}}\xout{\illegible{1}} $\underline{L}^{-1}$ \ comme un sous-\xout{\illegible{1}} module inversible de \ $\underline{E}$ locale
ment facteur direct dans $\underline{E}$, alors la correspondance entre $\hand{\xi}$ \ (i.e. $\underline{L}$ \hand{\unsure{$\xi$,}}

\mspage{\hand{4}}

\clearpage
% Nouveau EGA V (cote 25): archive cote 25, site page 140 (PDF page 141).
% typescript. Draft compared: none.
% AI-transcribed by Claude Opus 5.5 (claude-opus-5-5, adaptive thinking, effort high), Message Batches API, batch pilot (msgbatch_0189Nhf1xB615YWnztubtLfc), 2026-11-12; not yet checked by eye.
\scanpage{25/140}
\margin{\uline{nouveau EGA V}}

IV 16

Un nouveau n$^{\text{o}}$ : \uline{Ensemble singulier et supersingulier d'une fonction},
\uline{et critères différentiels}.

Ce n$^{\text{o}}$ sera utilisé au par.20 des sections hyperplanes, mais sa
place naturelle me semble bien au par. 16.

\uline{Définition} 1 Soient X un préschéma régulier, $\phi$ une section de $\underline{O}_X$ . Un
point $x \in X$ est dit \uline{zéro} \xout{\illegible{3}} \uline{singulier} pour $\phi$
si on a $\phi_x \in \underline{m}_x^2$, zéro \uline{supersingulier} si c'est un zéro singulier
et si de plus l'élément de $\underline{m}_x^2/\underline{m}_x^3 \simeq \mathrm{Sym}(\underline{m}_x/\underline{m}_x^2)$ qu'il définit,
interprété comme une forme quadratique sur le dual $\underline{t}_x$ de $\underline{m}_x/\underline{m}_x^2$ sur k(x),
est une forme \deleted{non} dégénérée.( Un zéro singulier qui n'est pas supersingu-%
lier est parfois appelé \uline{zéro singulier ordinaire}).

\uline{Remarques} 2 Si $x \in V(\phi)$\hand{,} alors x est un zéro \inserted{non} singulier de $\phi$ sss
$\phi_x \neq 0$ et x est un point \xout{\unsure{régxxx}} non singulier i.e. régulier de $V(\phi)$,
i.e. sss x est un point régulier de $V(\phi)$ et $V(\phi) \neq X$ au voisinage de
x . \deleted{On peut prouver que si x est singulier pour $\phi$ , il est singulier
ordinaire sss $\phi_x \neq 0$ et la multiplicité de $\underline{O}_{V(\phi),x}$ est égale à 2
(ou, comme on dit encore parfois, x est un \uline{point singulier ordinaire}
\uline{du préschéma} $V(\phi)$ ); donc un $x \in V(\phi)$ est non supersingulier \inserted{pour $\phi$} sss
$\phi_x \neq 0$ et la multiplicité de $O_{V(\phi),x}$ est $\leq 2$.}

\uline{Définition} 3. Soient X un préschéma lisse sur un corps k, $\phi$ une
section de $\underline{O}_X$ , $x \in V(\phi)$ . On dit que x est un zéro \uline{géométriquement}
\uline{singulier} \xout{\unsure{dexx}} (resp. \uline{géométriquement \unsure{supersingulier}} ) de $\phi$ rel à k,
si pour toute extension k' de k et tout point \hand{$\xi$} de X à valeurs dans
k\hand{,} localisé en x, la \xout{\illegible{1}} point correspondant x' de $X_{k'}$, est un zéro
\xout{\unsure{géom}} singulier (resp. supersingulier) de $\phi_{k'}$ .

\mspage{140}

\end{document}
